\section{Related Work}
\label{sec:related}

\paragraph{Spatial-domain detection.}
FaceForensics++~\cite{rossler2019faceforensics} established the first
large-scale benchmark for face manipulation detection, showing that an
XceptionNet fine-tuned on RGB frames achieves strong accuracy across five
manipulation types.
The success of ImageNet-pretrained CNNs on deepfake detection has since been
replicated across many architectures, suggesting that low-level texture
statistics from natural images transfer well to identifying synthetic
artefacts.

\paragraph{Frequency-domain analysis.}
Durall \etal~\cite{durall2019unmasking} showed that GAN images exhibit
systematically elevated high-frequency energy.
A simple classifier on the 1D azimuthal power spectrum separates real from
GAN-generated faces with reasonable accuracy.
Frank \etal~\cite{frank2020frequency} extended this with 2D DCT features,
achieving $>90\%$ detection rates for several GAN types, while also
documenting that frequency-based detectors degrade under JPEG compression.
Dzanic \etal~\cite{dzanic2020fourier} provided the theoretical basis:
the periodic ``checkerboard'' pattern in frequency space arises from uneven
kernel overlap in transposed-convolution upsampling layers.

\paragraph{Explainability.}
CAM~\cite{zhou2016cam} and Grad-CAM~\cite{selvaraju2017gradcam} are now
standard tools for interpreting CNN decisions.
In deepfake detection, heatmaps have consistently shown that
well-performing models attend to periocular and mouth regions — areas where
GAN synthesis tends to concentrate synthesis errors.
We apply Grad-CAM to validate that our models are not relying on spurious
background cues.
